\chapter{Background and Related Work}
\label{ch:related-work}

The central distinction in this review is between improving a tutor's response, estimating performance, and deciding which
additional evidence to collect. These tasks can share software and mathematical tools while requiring different evaluations.
The literature is compared through those decisions and outcomes, rather than through the number of agents or models used.

AI in education includes learner modelling, automated assessment and tools that support teaching. Reviews of the field warn
against judging an educational system solely by what its software can do. The educational purpose, the people involved and the
outcome being measured also matter \parencite{roll2016evolution,zawackirichter2019ai}. This report therefore keeps its
engineering demonstration separate from its prediction experiments.

The wider feedback literature reaches a related conclusion: the value of feedback depends on its timing, content and the
learner's opportunity to use it \parencite{hattie2007power,shute2008formative}. That does not make this a feedback-effectiveness
study. It explains why a tutor needs a defensible basis for deciding whether to seek more evidence before it acts.

\section{Why Evidence Decisions Matter Before Tutor Action}

An AI programming tutor must decide what to do after a learner's action. It may offer guidance, ask for more evidence, or move
on. The present dissertation studies the evidence decision that comes before that tutor action. It does not study how feedback
is worded or whether a particular explanation improves learning.

The educational case for intelligent tutoring systems is substantial, but it is not uniform. Meta-analyses report positive
average effects compared with conventional instruction, while also finding that outcomes depend on the comparison condition,
the assessment used and the match between what the tutor teaches and what the test measures
\parencite{ma2014intelligent,kulik2016effectiveness,vanlehn2011relative}. This makes the evaluation target important. A
system that produces plausible assistance during practice has not thereby shown that learners can solve a related problem
independently.

Programming education has a long history of automated assessment. Unit tests, reference answers and program analysis make
observations available quickly. Surveys show that the assessment target must be stated carefully, because dynamic tests,
static checks and feedback tools examine different properties of a program \parencite{alamutka2005assessment,ihantola2010assessment}.
They do not, by themselves, establish which programming idea a learner understands. Prior work on automated feedback also
shows why the assessment boundary matters: immediate task performance does not necessarily show whether a learner can solve a
later task independently \parencite{keuning2018systematic,cavalcanti2021automatic}. This dissertation uses that distinction
to define its setting. It does not evaluate feedback wording or pedagogical quality. A recent review of AI in programming
education likewise shows a varied set of applications and evaluation choices, rather than one settled measure of educational
value \parencite{manorat2025artificial}.

The need for support should also be stated carefully. Reviews describe fragile knowledge, incomplete mental models and recurring
misconceptions among novice programmers \parencite{robins2003programming,qian2018misconceptions,luxtonreilly2018introductory}.
Claims that introductory programming has an unusually high failure rate are less settled. An early international survey had
too few participating institutions for firm conclusions \parencite{bennedsen2007failure}. A later review of 161 courses found
an average pass rate of 67.7\%, with limited variation by language and over time \parencite{watson2014failure}. These findings
support studying programming evidence without treating every error as proof that programming education is failing.

LLMs make the surrounding tutor more fluent, but they also make weak evidence harder to notice. Reviews of LLM use in
education report limited replicability, low technology readiness and unresolved concerns about transparency, privacy and
beneficence \parencite{yan2024practical}. Experimental reviews also warn that improvements in assisted performance do not by
themselves establish durable learning, and call for stronger outcome measures and adequate power \parencite{deng2024chatgpt}.

This dissertation therefore focuses on a decision that sits before any tutor explanation: whether an additional executable check
should be requested and how much trust it should receive. The question is narrower than whether an LLM tutor improves learning.
It asks whether the system can make a more reliable prediction when available evidence may be incomplete, irrelevant or wrong.

\section{Learner Modelling and Programming Evidence}

Bayesian knowledge tracing separates observed performance from an inferred knowledge state. In the original programming-tutor
model, each coding rule is either learned or unlearned. A learned rule can produce an incorrect answer through a slip. An
unlearned rule can produce a correct answer through guessing. The model also includes initial knowledge and a probability of
learning during practice. Its original formulation assumes no forgetting
\parencite[Section 3, pp.~260--261, Figure 4]{corbett1994knowledge}.

Observation and learning are separate steps. If $q$ is the probability that a rule was already learned after accounting for
the observed answer, and $\tau$ is the learning-transition probability, the next estimate is
\begin{equation}
 p_{\mathrm{next}}=q+(1-q)\tau.
\end{equation}
This restates Corbett and Anderson's equation~1: the rule was already learned, or it was unlearned and is acquired during the
opportunity to practise \parencite[p.~261]{corbett1994knowledge}. Updating an estimate from evidence and modelling a change in
knowledge are therefore different operations.

That distinction constrains this dissertation. The external-model benchmark's tracker applies fixed score changes. The
bounded-update study evaluates belief updates against authored hidden states, while the probe-selection study starts with supplied probabilities and performs
one batch selection without a learning transition. Neither a rising demo score nor a successful simulated update establishes
human learning. The shared contribution of knowledge tracing is the explicit observation/state distinction, rather than evidence
that these simpler project components measure mastery.

Deep Knowledge Tracing uses recurrent neural networks to learn a representation of interaction history and predict future
performance. Its inputs encode the exercise and whether the answer was correct. The learned state can capture dependencies
across interactions without an explicit binary mastery variable for each skill
\parencite[Section 3]{piech2015dkt}. Fitting this representation requires learner records, and its state is less directly inspectable.
The present studies use transparent, authored probabilities for a different reason: they test how a prediction changes after
one additional coding check. They do not fit a learner model from the 23 authored cases.

Learner models also differ in what they allow to vary. Reviews cover probabilistic, cognitive and item-based approaches
\parencite{desmarais2012learner}. Individualised Bayesian knowledge tracing has shown that learner-specific rates can improve
prediction for unseen students. In that evaluation, the model could not use fitted personal parameters for the test learners.
The authors attribute the gain to better shared skill parameters after accounting for learner variation during fitting
\parencite[Section 4]{yudelson2013individualized}. The project does not estimate such rates. Its historical
study uses simple predictors with learner-separated folds, while its simulators hold their learner assumptions fixed. This
limits any claim that the resulting probabilities describe an individual learner's knowledge.

TIKTOC makes the programming observation more detailed: it jointly predicts student code and whether that code passes
individual tests. Its training objective combines code-generation and test-outcome losses
\parencite[Section 4.3, equation 4]{duan2025tiktoc}. This is directly relevant to the original proposal, since using executable
test information to model programming performance is already established work. The present dissertation cannot claim that
idea as its novelty.

The distinction lies in the decision being studied. TIKTOC learns predictions from interaction histories and test-level
information. The probe-selection study allocates a fixed number of additional observations under uncertain channel reliability. TIKTOC's reported
model comparison does not evaluate that allocation problem, and this project does not reproduce its learned predictor.
Consequently, the results cannot establish superiority over TIKTOC.

TIKTOC augments Java CodeWorkout data with 305 tests across 17 problems. The reported
subset contains 3,714 submissions from 246 students \parencite[Section 3.2, Table 2]{duan2025tiktoc}. This dissertation's study of
Python CSEDM records selects uncertain first-attempt predictions for hypothetical review. It neither reconstructs
TIKTOC's test-level targets nor evaluates optional executable checks. Its larger observational sample therefore provides a
separate finding, rather than external validation of the probe-selection rule.

TRAVER combines knowledge tracing with a trained turn-level verifier that ranks candidate tutor utterances
\parencite[Sections 3.1--3.2]{wang2025traver}. Its DICT protocol evaluates coding tutoring through simulated students and code
generation tests. The selected object is therefore a tutor response, whereas this dissertation's probe-selection study selects cases for additional evidence.
The outcomes also differ. DICT assesses task completion in a tutoring process. This dissertation evaluates prediction under
fixed evidence conditions.

This distinction prevents a misleading comparison. The project does not outperform TRAVER, and its frozen results cannot be
placed on the same performance scale. TRAVER establishes that combining tracing and verification is already part of the
coding-tutor literature. The narrower issue investigated here is how evidence handling and selection behave when the assumed
evidence channel is wrong. That is a difference in the experiment, not proof of a globally new architecture.

\subsection{What an Uncertain Observation Means}

A confidence score needs an interpretation before it can support a probability update. Chan and Darwiche distinguish two
ways to specify uncertain evidence: Jeffrey's rule takes prescribed revised probabilities for a set of events, whereas virtual evidence
specifies evidence strength through likelihood ratios. They explain how to translate between these descriptions when the
prior is known \parencite[pp. 69--74, 77--78]{chan2003uncertain}. A statement about confidence does not, by itself, say which
description is intended.

Consider a success probability of $0.2$, giving prior odds of $1:4$. Evidence with a likelihood ratio of $4$ changes those
odds to $1:1$, so the revised probability is $0.5$. Requiring a revised probability of $0.8$ instead gives odds of $4:1$.
Reaching that value from the same prior would require a likelihood ratio of $16$. This illustrative calculation shows why
evidence strength and revised belief cannot be substituted for each other.

The distinction matters for the bounded-update study. Its channel-aware trackers use authored probabilities of an observation given each simulated
state. They retain the input confidence field but do not use it in the update. The legacy comparison applies a separate
fractional likelihood transformation to confidence. That transformation should not be presented as an implementation of
Jeffrey's rule or as evidence that confidence is calibrated. The update equations in \cref{ch:method} specify what the
channel-aware rules actually do. Clipping then modifies their contribution, so the resulting belief need not be the posterior
under the original observation model. These are modelling choices whose consequences require evaluation.

\section{Active Evidence Acquisition}

Active-learning research provides the broad idea of asking for information where it may be most useful. Early work used model
uncertainty to select training examples, and later surveys distinguish that decision from prediction itself
\parencite{lewis1994sequential,cohn1996active,settles2009active}. Those methods commonly acquire labels to improve a model
for future use. The current study instead allocates extra checks to cases that are already being predicted. This difference
matters because an attractive training-data acquisition rule need not improve a saved case-level prediction.

Dynamic LENS maintains a distribution over learner-state estimates, combining learned forecasts with encodings of observed
responses \parencite[Section 2]{christie2024dynamiclens}. Its evaluation examines prediction and how uncertainty changes as
history is included. The discussion proposes selecting questions by expected information gain, while noting that efficient
approximation needs further investigation \parencite[Sections 3--4]{christie2024dynamiclens}. It supports the motivation for
uncertainty-aware assessment. It is not a reported fixed-budget acquisition trial against which this dissertation's results
can be directly compared.

The project's uncertainty-only baseline uses a simpler quantity: closeness of a single predicted success probability to
one half. This does not distinguish uncertainty about an estimated probability from variation that remains even when that
probability is known. For example, a known fair coin has outcome probability one half, but further flips need not reveal
anything about its already-known bias. A poorly estimated bias may also have mean one half, with a different reason for
collecting data. The example illustrates why equal prediction scores can conceal different opportunities for improvement.

Calibration asks a separate question: whether predicted probabilities match observed frequencies. Modern classifiers can be
accurate while poorly calibrated \parencite{guo2017calibration}, and uncertainty methods can change their behaviour when the
data distribution shifts \parencite{ovadia2019uncertainty}. These findings motivate the held-out fault settings in this
dissertation. They do not show that the authored settings match a classroom, or that the project's probabilities are calibrated
outside their declared simulator.

The historical learner-record study tests whether a simple uncertainty ranking identifies errors in saved predictions. A positive result does not establish
that those errors are reducible through review. The probe-selection study additionally models the evidence channel, but its score must still account
for how an observation affects the executed decision. Neither study implements Dynamic LENS or separates these sources of
uncertainty through a learned distribution over learner states.

Prediction-oriented active learning makes a related distinction between learning model parameters and improving predictions
on inputs of interest. \textcite[Sections 2.1 and 4]{bickfordsmith2023prediction} introduce expected predictive information gain
(EPIG). It measures how informative an acquired label is about a target label, averaged over a specified target input
distribution. Their active-learning procedure acquires labels for training and updates the predictive model.

The target distribution matters: information about a remote or irrelevant part of the input space may contribute little to
the predictions that the system needs to make. EPIG expresses relevance through expected predictive information, rather than
through a label such as ``related task''. Its connection to expected generalisation error uses cross-entropy under the stated
model \parencite[Section 4, equation 5]{bickfordsmith2023prediction}. It is not a guarantee that every acquired label reduces
realised classification error.

This provides a useful comparison for the dissertation's design. The probe-selection study neither implements EPIG nor retrains a predictive model
after acquiring labels. Its checks update separate case beliefs under fixed calibration, and its endpoint is binary
classification error. The relevant lesson is to specify which future prediction an observation should help, then evaluate
the actual decision rule. A shared prediction-oriented motivation does not make the clipped-belief score equivalent to EPIG,
or allow its negative result to stand as a comparison against that method.

An uncertain prediction is not necessarily one that an extra observation can improve. For the binary decision studied here,
uncertainty-only selection ranks probabilities nearest to one half. A decision-value calculation instead needs a model of the
possible observations and the action that would follow each one. Expected-loss minimisation is established decision theory
\parencite[Section 5.7]{murphy2012machine}. The dissertation's implementation must be judged against that definition.

This distinction matters for the completed experiment. The probe-selection study uses a change in clipped-belief scores to rank candidates, while
the measured outcome is classification error. These quantities need not agree, as the counterexample in \cref{ch:method}
shows. Calling the score ordinary value of information would therefore obscure a methodological difference. Its lower
calibration quantile also represents caution within the assumed model. It gives no protection against failure mechanisms
that the model omits.

Tang and colleagues examine the related problem of Bayesian active learning under model misspecification. Their analysis
considers how selected training data can differ from the target distribution and how updates can amplify or reduce error.
R-IDeA incorporates representativeness, information and error-reduction considerations into acquisition
\parencite[Sections 3--4]{tang2026robust}.

The theoretical setting is specific. Theorem 1 concerns squared-error generalisation for an empirical risk minimiser in a finite
model class, with bounded functions and outcomes. The bound includes a training-to-test density-ratio factor. These conditions
matter. Remark 3 distinguishes that empirical-risk-minimisation assumption from the paper's Bayesian experiments, which predict
using a distribution over models. The theorem informs acquisition design, but should not be read as a guarantee for any rule
described as Bayesian or reliability-aware. This comparison checks the stated assumptions, rather than independently verifying
the complete appendix proof.

The probe-selection study does not implement R-IDeA or test its theoretical bounds. It performs a single batch allocation, keeps calibration fixed
and evaluates deliberately altered observation channels. Its lower quantile and bounded update are different mechanisms.
It neither fits the theorem's empirical risk minimiser nor establishes its bound for the binary classification endpoint.
The connection is the possibility that the same wrong model can distort both evidence selection and the interpretation of
what is collected. This motivates separate evaluation environments. It does not validate the project's chosen fault rates
or establish that its heuristic is robust.

The completed comparison consequently supports a limited methodological lesson. A cautious-looking acquisition score must
still be checked against the decision that is executed, and its empirical performance must be measured when assumptions
change. The frozen negative result concerns this implementation. It cannot establish that Bayesian acquisition methods in
general fail, or that R-IDeA would perform better on these cases without a separate comparison.

\subsection{Selecting Tests When Their Outcomes Are Noisy}

Noisy test selection already has a substantial mathematical treatment. Chen and colleagues' ECED method selects tests
sequentially to identify a hidden target, using a surrogate objective and an error-based stopping rule
\parencite[Sections 2--3]{chen2017noisy}. The model specifies a known distribution over root causes and known observation
parameters. Tests are conditionally independent given a root cause, although they may remain dependent given the target.
Each test has unit cost, and persistent noise means that repeating it returns the same outcome.

Theorem 1 bounds the number of tests under binary outcomes and test-specific noise rates independent of the root cause.
Its reference is an optimal policy achieving a stricter error target \parencite[Section 3, Theorem 1]{chen2017noisy}.
This is not a guarantee for arbitrary correlated corruption or an unknown, misspecified observation model. The comparison
here checks the theorem's statement and assumptions. It does not independently verify the supplementary proof.

This sharpens the boundary of the dissertation's contribution. The probe-selection study estimates channel reliability, selects a fixed batch and
clips belief updates. It neither implements ECED nor proves a link between its ranking score and optimal test selection.
Although both methods use a surrogate score, ECED's stated theoretical connection cannot be transferred to the project's heuristic by analogy.
Conversely, the heuristic's negative result says nothing about how ECED would perform in an appropriately defined comparison.

The useful distinction is therefore the particular decision rule and evaluation, rather than a claim to introduce noisy
evidence acquisition. The project examines a limited batch heuristic under specified faults and identifies where its score
and executed decision diverge. A sequential method with a stopping rule would answer a different question and require a
separate experimental design.

Jia and colleagues study noisy decision trees with a different objective: identify the hidden hypothesis while minimising
expected test count \parencite[Sections 1--2]{jia2024noisy}. Their binary model supplies a hypothesis prior and the positions
of noisy entries in a test matrix. Outcomes at those entries are uniform and conditionally independent given the hypothesis.
Repeating a test returns its previous outcome. The main identification assumption ensures that deterministic tests can
distinguish any pair of hypotheses. Section 7 considers relaxations.

The authors explicitly distinguish exact identification from the permitted error in ECED
\parencite[Section 1.2]{jia2024noisy}. This places both methods within established noisy-test selection, with different
assumptions and objectives. The probe-selection study instead estimates potentially wrong channel reliability and measures classification error
after a fixed allocation. It does not implement either algorithm or inherit their approximation guarantees. These differences
explain why a direct superiority claim would be inappropriate. They do not establish novelty by themselves.

\subsection{Limiting the Influence of Misleading Observations}

Robust filtering addresses what happens after an observation is acquired. Duran-Martin and colleagues propose the weighted
observation likelihood filter (WoLF), which uses generalised Bayesian updating to reduce the influence of outlying measurements
\parencite{duranmartin2024outlier}. Their approach weights the observation log-likelihood and derives closed-form updates for
Gaussian filtering, with extensions to nonlinear settings. The comparison here uses Sections 3.1--3.4 of the
\href{https://utstat.utoronto.ca/~alexander/wolf-icml.pdf}{co-author-hosted manuscript}.

Their robustness statement has a specific meaning. Section 3.4, Theorem 3.2 considers linear Gaussian models or their linearised
approximation and imposes conditions on the weighting function. It bounds the divergence between posteriors obtained from an
original observation and an arbitrarily contaminated replacement. That is a statement about posterior sensitivity under those
conditions. It should not be substituted for a claim about classification accuracy or learning.

The bounded-update study takes a simpler route. It caps the log-likelihood-ratio contribution to a binary belief update and applies fixed channel
trust. The cap bounds one observation's contribution to log-odds. It neither implements WoLF nor establishes its posterior
influence theorem. Repeated misleading updates can still accumulate, and clipping can discard useful evidence as well as limit
harmful evidence. The project's adverse-condition and clean-condition comparisons test that trade-off under authored rules.

Together, robust filtering and active acquisition suggest two separate design questions: how strongly to update after a check,
and which check to request beforehand. This dissertation evaluates both through distinct studies, but their combination alone
does not establish a new robustness theorem. Its claim rests on the measured behaviour and stated limitations of the implemented
rules.

\subsection{Learning How Much Weight to Give Existing Evidence}

Zhang and colleagues describe a reinforcement-learning gate that combines behavioural and dialogue-derived prediction
streams. The gate assigns a convex weight to their logits. Its training objective includes predictive utility, a cohort
disparity penalty and entropy regularisation \parencite[Sections 3.2--3.4]{zhang2026fusion}. This is a close precedent for
adaptive evidence weighting in educational prediction.

The action differs from probe allocation. The gate weights features already available to the predictor. This dissertation's selector chooses which
cases receive another observation, using a fixed ranking rule rather than a policy learned through reinforcement learning.
The bounded-update study's channel trust is also fixed. The dissertation therefore does not reproduce the learned gate or compare its performance
with the proposed method.

The source's evaluation descriptions require caution. Section 3.1 groups learners by total practice volume, whereas
Section 4.1 describes within-learner progress using $t/T$ \parencite{zhang2026fusion}. Those groupings answer different
questions. The architectural comparison remains useful, but this review does not adopt the reported fairness gains or
resolve that ambiguity. More broadly, evidence weighting, evidence acquisition and changes in learning each require their
own outcome and comparison. Sharing an adaptive-learning application does not make them interchangeable.

\subsection{Acquisition Cost and Fixed Allocations}

Melville and colleagues study the acquisition of missing feature values for classifier training. Their score averages the
estimated improvement in classifier accuracy over possible feature values, divided by acquisition cost
\parencite[Section 2, equations 1--2]{melville2005acquisition}. The classifier is retrained as values are acquired.
This establishes a direct precedent for choosing observations according to expected benefit and cost. The paper's main
experiments assume equal acquisition costs. Their results should not be described as measurements of monetary savings.

The probe-selection study differs in both timing and objective. It allocates checks to evaluation cases, retains its fitted calibration model and
updates individual predictions. Every attempted probe counts as one unit, including a missing result. The primary policies
receive the same allocation, so differences in their error rates compare which checks they select. They do not show that
one policy uses fewer calls than another at that allocation.

The secondary budget curve compares allocations, but does not price their consequences. A deployment would need to account
for model requests, execution, retries, elapsed time and the burden of answering an extra question. Different checks could
have different costs. No such cost estimate enters the frozen selector, and the simulator makes no external requests.
The defensible description is therefore fixed-budget evidence selection. The resource records describe observed workload
in the external-model experiment. They do not establish realised savings from a deployed tutor.

Selective prediction provides a nearby, but different, decision. A selective classifier can refuse to make a prediction when
its expected risk is too high, trading coverage against error \parencite{geifman2017selective}. The acquisition experiment
cannot refuse: it must allocate exactly 20 checks. A low-scoring case therefore loses a check to another case. It does not
save an interaction. This is why the study reports error at a fixed allocation rather than a risk--coverage guarantee.

The same distinction now appears directly in knowledge tracing. Mitton and colleagues use model uncertainty to defer a fixed
share of predictions for teacher review across three neural knowledge-tracing models \parencite{mitton2026defer}. Their work
supports uncertainty as a review signal. It does not test whether an extra coding check improves the deferred prediction. The
CSEDM companion study addresses the first question with simpler models. The acquisition study addresses the second only inside
its declared simulator.

\section{Evaluation and Simulator Validity}

The evaluation must match the decision being studied. A teaching intervention asks whether assistance changes learning.
A prediction study asks whether observations anticipate an outcome. An acquisition study asks which additional observations
improve that prediction within a resource limit. Evidence for one question can motivate another, but the connection requires
an experiment. This section uses that distinction to assess the available evaluation approaches.

Tutor-response quality, prediction accuracy, and later independent performance are different outcomes. Dr.Academy evaluates
the quality of educational questions \parencite{chen2024dracademy}. MRBench evaluates tutor replies to mistakes and confusion
across eight pedagogical dimensions \parencite[Sections 3--4]{maurya2025tutorevaluation}. MathTutorBench combines several
pedagogical evaluation tasks \parencite{macina2025mathtutorbench}. The BEA shared task also found that automatically judging
several tutoring dimensions remains difficult, even with human-labelled material \parencite{kochmar2025bea}. These criteria
expose aspects of tutoring that a code pass rate cannot measure. A correct answer can still give away the solution or leave the
learner without a useful next step.

Evaluation software can make these distinctions easier to inspect. AITutor-EvalKit provides tutor assessment, model inspection
and data visualisation in one application \parencite{naeem2026evalkit}. Earlier conversational systems such as AutoTutor also
combined dialogue, feedback and scaffolding, with learning evaluated across several domains \parencite{nye2014autotutor}.
These systems show the breadth of a tutor evaluation. The project's interface demonstrates its evidence flow, while its reported
experiments measure prediction and evidence handling.

The v1 public-answer ratings serve a different purpose. They categorise the evaluation model's visible answer as input to a
prediction rule. They do not assess the tutor's pedagogical quality. Reviewer agreement therefore cannot substitute for an
MRBench-style tutor evaluation. Conversely, favourable ratings of the demonstration would not establish that its predictions
are accurate or that learners improve. This project preserves separate claims for the interface and the research results.

LongTutor separates evidence acquisition, state diagnosis and teaching action in an evaluation grounded in learning histories
\parencite[Section 3.1]{li2026longtutor}. Its evidence task asks models to extract historical information, combine records and
recognise false premises. This differs from acquiring a new executable observation. The distinction is useful here: finding
relevant history, predicting a later outcome and choosing guidance each require their own assessment. The historical learner-record study's
features do not amount to LongTutor's diagnosis task, and the project's interface has not been evaluated on its teaching criteria.

Two recent alternatives sharpen this boundary further. Koutcheme and colleagues use program repair as a proxy when evaluating
LLM feedback for programming education \parencite{koutcheme2024repair}. The proxy may be practical, but it does not measure
whether the learner understands the explanation. EducationQ evaluates teaching through multi-agent simulated dialogues
\parencite{shi2025educationq}. Its design shows one way to study teaching behaviour, while also reinforcing the need to state
what a simulated interaction can and cannot establish. Neither approach evaluates the case-level evidence decision used here.

\subsection{A Simulated Response Is Not a Learner Outcome}

The response to feedback needs its own validation. \textcite[Sections 2.1--2.3]{do2026misconception} compare a model's revision of
an initially wrong answer under targeted, misaligned and generic feedback. Their Selective Flip Score subtracts the mean
correction rate under the two controls from the correction rate under targeted feedback:
\begin{equation}
 \operatorname{SFS}=F_T-\tfrac{1}{2}(F_M+F_G).
\end{equation}
Here a correction means that the revised answer is correct. Across their seven prompted models, the authors report similarly
high correction rates across feedback conditions and near-zero SFS. Their preprint examines behaviour that answer similarity
alone would miss: whether a model responds selectively to relevant guidance.

The score requires careful interpretation. Zero also occurs when all three correction rates are zero. On its own it cannot
identify a model abandoning a misconception and solving the problem again. The component rates and response analysis are
needed to distinguish that behaviour from never correcting an answer. This follows from the score's definition, and limits
the mechanism that can be inferred from an aggregate difference.

The distinction matters for v1. Its controls change the evidence supplied to a fixed prediction rule, rather than the feedback
given to a model revising the same answer. Equal predictions after related and unrelated passing checks can arise directly from
the tracker's category-based update. They do not establish the simulator behaviour studied by Do et al. This dissertation
therefore uses their contrastive design as related methodological context. The v1 study is neither an SFS replication nor a validation
of an external model as a student.

StudentSim treats simulator construction as a separate training problem. It distinguishes matching an individual learner's
responses from changing a response under guidance, and uses pooled training followed by individual specialisation
\parencite[Sections 3--4]{yang2026studentsim}. Its evaluation spans chess, English writing and mathematics. Its tutor-training
demonstration concerns chess, with expert ratings of generated guidance. These ratings should not be read as a measured
learning gain for programming students. The source is a recent preprint, and its scope remains important when drawing comparisons.

This dissertation has neither fitted its external model to learner records nor validated its reactions to teaching. A plausible
wrong answer would not establish either property. Calling that model an evaluation target keeps the measured quantity explicit:
the pipeline predicts whether that model's separate code response passes the authored checks. The result is about the model
under those prompts, not an estimate of how students would learn from the tutor.

The mathematical simulators serve a different purpose. Their hidden states and observation faults are specified so that
updates and selection rules can be tested under controlled changes. Independent review can challenge their internal logic.
It cannot establish that the chosen rates describe learners. The CSEDM records add observations from people, but contain no
matching optional-probe intervention. Combining these studies therefore does not close the gap to a human tutoring evaluation.

\subsection{What the Prediction Metrics Measure}

Classification error assesses the chosen binary answer. It discards the distance from the decision threshold: probabilities
of 0.51 and 0.95 both predict success at a threshold of 0.5. Brier score and log loss also assess the probability assigned to
the observed outcome. Gneiting and Raftery describe their strict propriety: in expectation, reporting the true probability
uniquely optimises the score \parencite[Section 3, Examples 1--2]{gneiting2007strictly}. The paper uses rewards to maximise.
This dissertation reports losses to minimise.

For example, when a task succeeds, moving its predicted probability from 0.6 to 0.8 reduces binary Brier loss from 0.16 to
0.04 while leaving classification error unchanged. This arithmetic illustrates why the tracker and acquisition studies can
reach different conclusions. Their primary outcomes assess different properties. A lower observed Brier loss also does not
by itself prove calibration, or establish that requesting another check was useful. The report interprets each study using
its specified primary outcome and treats the remaining metrics as supporting evidence.

\subsection{Choosing an Evaluation Approach}

A student intervention would provide the most direct evidence for a claim about teaching. It would require comparable
starting conditions and an assessment completed without the tutor. The assisted and unassisted outcomes in
\textcite{bastani2025guardrails} show why that choice matters. Such a study would assess the combined effects of the interface,
guidance and learner response. Identifying the effect of probe selection itself would require a comparison that changes that
selection while controlling the other components. The present experiments do not supply this causal comparison.

Historical learner records offer observed programming performance and allow predictions to be tested on learners excluded
from fitting. They also constrain the questions that can be answered. CSEDM does not record the optional checks required by
the proposed acquisition policy, so it cannot reveal what would have happened after a different probing decision. Its role
here is to test whether uncertainty identifies errors worth investigating. Error detection alone leaves the value of a new
check unresolved.

Executable benchmarks make correctness reproducible under specified tests. Their weakness is that the tests define the
measured outcome: incomplete coverage can reward a response that violates the task. Holding the outcome back protects the
prediction procedure from seeing the answer, but it cannot repair a defective assessment. This motivates the correction
audit in the external-model study and limits its claim to the authored tasks and evaluation model.

Controlled simulation permits each selector to face the same hidden outcomes and potential observations, including checks
it chooses to skip. It therefore supports a direct comparison of acquisition rules. The price is dependence on the chosen
data-generating assumptions. Separate calibration and evaluation settings challenge those assumptions. They do not establish
how frequently the faults occur in education. The dissertation reports performance within those settings and uses the
external-model and historical studies to examine different parts of the evidence problem.

\section{Research Gap and Study Rationale}

The project tests a limited connection between evidence and prediction. Its external-model experiment preserves a record of
predictions made before a separate outcome was revealed. The tracker study tests bounded updates under authored faults, and
the acquisition study compares a fixed-budget heuristic with simpler selectors. The historical-data companion concerns
which prediction errors uncertainty can identify. None of these comparisons evaluates the quality of a tutoring dialogue.

This leaves a clear boundary for the contribution. Tracing, verification and pedagogical evaluation already have dedicated
methods. The present work contributes an inspectable experimental record and specific findings about its evidence-handling
rules, including their failures.

Table~\ref{tab:nearest-work} brings together the nearest comparisons. They share individual components with this dissertation,
but study different decisions or outcomes. Their published results cannot therefore be placed on the same performance scale.

\begin{table}[htbp]
  \centering
  \small
  \begin{tabularx}{\textwidth}{@{}>{\raggedright\arraybackslash}p{0.15\textwidth}>{\raggedright\arraybackslash}p{0.21\textwidth}>{\raggedright\arraybackslash}p{0.25\textwidth}>{\raggedright\arraybackslash}X@{}}
    \toprule
    Work & Decision studied & Evidence and reliability & Relation to this dissertation \\
    \midrule
    TIKTOC \parencite{duan2025tiktoc}
      & Predict student code and individual test outcomes from interaction history.
      & Learns from code, histories and test-level outcomes.
      & Establishes test-informed programming prediction. It does not allocate optional checks. \\
    \addlinespace
    TRAVER and DICT \parencite{wang2025traver}
      & Rank candidate tutor responses in simulated coding dialogues.
      & Uses a trained verifier and code-generation tests.
      & Selects tutor responses and measures task completion, rather than selecting evidence for a saved prediction. \\
    \addlinespace
    Dynamic LENS \parencite{christie2024dynamiclens}
      & Estimate learner state and uncertainty, with active question selection proposed.
      & Combines learned forecasts with encoded responses.
      & Motivates uncertainty-aware assessment, without reporting this fixed-budget check-allocation comparison. \\
    \addlinespace
    EPIG \parencite{bickfordsmith2023prediction}
      & Acquire training labels that inform predictions on target inputs.
      & Uses predictive information under a Bayesian model.
      & Updates a predictive model. The present study updates individual case beliefs and measures classification error. \\
    \addlinespace
    ECED \parencite{chen2017noisy}
      & Select tests sequentially to identify a hidden target and stop at an error bound.
      & Assumes a known observation model with specified noise.
      & Supplies noisy-test theory. Its guarantees do not apply to the project's fixed-batch heuristic or misspecified channels. \\
    \addlinespace
    R-IDeA \parencite{tang2026robust}
      & Acquire training data under model misspecification.
      & Combines representativeness, information and error-reduction criteria.
      & Motivates testing under changed assumptions. The project neither implements its rule nor inherits its bound. \\
    \addlinespace
    Uncertainty deferral \parencite{mitton2026defer}
      & Send a fixed share of knowledge-tracing predictions for teacher review.
      & Ranks model predictions by uncertainty.
      & Closest to the CSEDM error-triage question. It does not test whether another coding check improves the prediction. \\
    \addlinespace
    This dissertation
      & Allocate 20 optional checks among 40 simulated cases, then update their predictions.
      & Estimates channel reliability, bounds updates and tests altered evidence settings.
      & Adds an auditable prediction-before-outcome evaluation, but its acquisition score is not expected classification-error reduction. \\
    \bottomrule
  \end{tabularx}
  \caption{Nearest technical comparisons. The methods collect different evidence, make different decisions and measure different
  outcomes. The table identifies prior components and the narrower evaluation performed in this dissertation. It does not imply
  a direct benchmark comparison or a claim of superiority.}
  \label{tab:nearest-work}
\end{table}

Together, these studies rule out a broad claim that combining tests, uncertainty and selection is new. They leave a narrower
application question: how much should a programming prediction change after an extra check, and which checks help when their
assumed reliability is wrong? The dissertation studies that question through committed predictions, bounded updates, fixed-budget
selection and a separate check of uncertainty-based error triage on historical learner records.

The narrow contribution is the evaluation of these specified rules and the analysis of where their assumptions fail.
The review found no directly equivalent evaluation across these particular checks and fault settings, but this does not
establish global priority or superiority over the methods reviewed. In particular, the implemented acquisition score differs
from expected decision value, and its comparison cannot answer whether ordinary value-of-information selection would succeed.
That distinction is developed mathematically in \cref{ch:method} and carried into the interpretation of the results.
